\documentclass[10pt,a4paper]{amsart}
\usepackage[margin=19mm]{geometry}
\usepackage{amsmath,amssymb,mathtools}
\usepackage[scr=rsfs,cal=euler]{mathalfa}
\usepackage{xfrac}
\usepackage[unicode,hidelinks,bookmarks=false]{hyperref}
\newcommand{\ie}{i.e.\ }
\newcommand{\cf}{cf.\ }
\newcommand{\resp}{respectively\ }
\newcommand{\Subsection}{\subsection}
\providecommand{\nobreakdash}{\nobreak\hbox{-}}
\setlength{\emergencystretch}{2em}
\multlinegap0pt
\usepackage{bm}
\usepackage{bbm}
\usepackage{dsfont}
\usepackage{xspace}
\usepackage{cancel}
\usepackage{mathtools}
\usepackage{amsthm}
\makeatletter
\@ifundefined{thm}{\newtheorem{thm}{Thm}}{}
\@ifundefined{lem}{\newtheorem{lem}[thm]{Lemma}}{}
\makeatother
\DeclareMathOperator{\modulo}{mod}
\newcommand{\Z}{\mathbb{Z}}
\title[Chebotarev geodesic theorem: split case]{Chebotarev geodesic theorem: split case}
\author{Alberto Acosta Reche}
\date{Source: arXiv:2606.25903v1 (CC BY 4.0)}
\begin{document}
\maketitle

\noindent\emph{Source and licence.} Chebotarev geodesic theorem: split case. Version arXiv:2606.25903v1 (CC BY 4.0), distributed under the Creative Commons Attribution 4.0 International licence, as recorded in the arXiv licence metadata for this exact version and shown on the abstract page https://arxiv.org/abs/2606.25903v1. Full rights notice: licenses/arxiv-2606.25903-CC-BY-4.0.txt.\par\smallskip

\noindent\textit{Editorial scope.} Counted unit: the Lem whose source label is \texttt{lem:weilbound} in the article (source0.tex lines 883--888, source label \texttt{lem:weilbound}), delivered together with the article's own complete original proof of it (source0.tex lines 890--953, one \texttt{proof} environment of 64 lines). No mathematics is written, repaired or re-derived: statement and proof are the article's own text line for line. Required context retained uncounted: none. Every definition, display and label of those lines is retained unchanged. Omitted from this package and not redistributed: 1--882, 889--889, 954--3177. Nothing inside a retained excerpt is omitted or altered: every retained body is delivered exactly as the article's lines are written, so no omission receipt of any kind accompanies this package. Citations: the retained excerpts carry no citation command at all, so no bibliography entry of the article is transcribed and no reference is redistributed. Reference adaptation: none was needed and none was made; every reference inside the delivered unit resolves inside it, so no unresolved reference or citation is produced. Printed slips kept literal: no printed word, symbol, hypothesis or number of the counted statement or of its proof was altered. Layout and class: the article's own class is \texttt{article} with packages including amsmath, amsthm, amssymb, mathtools, leftindex, tikz-cd, verbatim, bbm, enumitem, parskip, geometry, hyperref, biblatex; the wrapper is the lane's pinned \texttt{amsart} editorial wrapper at 10pt on A4, which restates the theorem-like environments the retained text needs and adds the article's own macro definitions that the retained text needs (\texttt{\textbackslash Z}, \texttt{\textbackslash modulo}), with extra wrapper packages (bm, bbm, dsfont, xspace, cancel, mathtools). The delivered numbering is the unit's own. Assets: no retained span contains a figure environment, a graphics-inclusion command, a PDF-page inclusion, a table or a TikZ picture; no image file is copied into this package, the package declares no asset dependency and no image file is redistributed with it. Licence: version arXiv:2606.25903v1 of this article is distributed under the Creative Commons Attribution 4.0 International licence, as recorded in the arXiv licence metadata for this exact version and shown on the abstract page https://arxiv.org/abs/2606.25903v1. A full rights notice is delivered at licenses/arxiv-2606.25903-CC-BY-4.0.txt. Characters: every retained excerpt is pure ASCII, so the delivered engine, XeLaTeX with its default Latin Modern text font, reports no missing character.
\begin{lem}\label{lem:weilbound}
If $m, n \neq 0$, then for any $\varepsilon > 0$ we have
\begin{equation*}
 S_{(\Gamma, \chi)}(m, n; c) \ll_{q, \varepsilon} \gcd(qm, qn, c)^{1/2}c^{1/2 + \varepsilon}.
\end{equation*}
\end{lem}

\begin{proof}
Of course, we can assume that $m, n \in \ell^{-1}\Z - \alpha/\ell$. We let $m' := qm, n' := qn$, which are nonzero integers. After a small manipulation we find that 
\begin{equation}\label{auxeq:kloosterman:1}
    S_{(\Gamma, \chi)}(m, n; c)  = \frac{\ell^2}{q^2}\sum_{\tilde{\gamma} \in \tilde{\Gamma}}\overline{\chi(\tilde{\gamma})}\sum_{\substack{a, d \text{ mod }cq\\
    ad = 1 \modulo c\\
    \begin{psmallmatrix}
        a & (ad-1)/c\\
        c & d
    \end{psmallmatrix} = \tilde{\gamma} \text{ mod }q}} e\left(\frac{m'a + n'd}{cq}\right).
\end{equation}
We will estimate each inner sum separately. We may assume that $c = c(\widetilde{\gamma}) \modulo q$, since otherwise the inner sum vanishes. We can write $q = \prod_{i= 1}^t p_i^{r_i}$, where $t\geq 1$, and where each $p_i$ is prime and each is $r_i$ positive for $1 \leq i \leq t$. Write also $c = c_0 \prod_{i = 1}^t p_i^{s_i - r_i}$, where $(c_0, q) = 1$ and the product ranges over the primes dividing $q$. We have
$cq = c_0 \prod_i p_i^{s_i}$. We define
\begin{equation*}
    c_0' = c_0, \quad c'_i  = p_i^{s_i} \text{ for }1\leq i \leq t,
\end{equation*}
so that 
\begin{equation*}
    cq = \prod_{i=0}^t c_i', \quad (c_i', c_j') = 1 \text{ for } i\not = j.
\end{equation*}
We define also 
\begin{equation*}
    R_i = \prod_{j \neq i} c_j',
\end{equation*}
and let $\overline{R_i}$ be any integer which is the inverse of $R_i$ when considered modulo $c_i'$. From the Chinese remainder theorem we see that each inner sum in \eqref{auxeq:kloosterman:1} factors as 
\begin{equation*}
    \sum_{\substack{a, d \text{ mod }cq\\
    ad = 1 \modulo c\\
    \begin{psmallmatrix}
        a & (ad-1)/c\\
        c & d
    \end{psmallmatrix} = \tilde{\gamma} \text{ mod }q}} e\left(\frac{m'a + n'd}{cq}\right) = S^{\tilde{\gamma}}_0(m'\overline{R_0}, n' \overline{R_0}; c_0) \times \prod_{i = 1}^t S^{\tilde{\gamma}}_i(m'\overline{R_i}, n' \overline{R_i}; c; p_i^{s_i}),
\end{equation*}
where 
\begin{equation*}
S^{\tilde{\gamma}}_0(m'\overline{R_0}, n' \overline{R_0}; c_0) = \sum_{\substack{a, d \text{ mod }c_0\\
ad = 1 \text{ mod }c_0}} e\left(\frac{m'\overline{R_i}a + n'\overline{R_i}d}{c_0}\right)
\end{equation*} 
is the classical Kloosterman sum and for $1\leq i \leq t$ we have
\begin{equation*}
S^{\tilde{\gamma}}_i(m'\overline{R_i}, n' \overline{R_i}; c; p_i^{s_i}) = \sum_{\substack{a, d \text{ mod }p_i^{s_i}\\
(a, d) = (a(\tilde{\gamma}), d(\tilde{\gamma})) \text{ mod }p_i^{r_i}\\
ad = 1 + b(\tilde{\gamma})c \text{ mod }p_i^{s_i}}} e\left(\frac{m' \overline{R_i}a + n' \overline{R_i}d}{p_i^{s_i}}\right).
\end{equation*}
The Weil bound gives
\begin{equation*}
S^{\tilde{\gamma}}_0(m'\overline{R_0}, n' \overline{R_0}; c_0)     \ll_\varepsilon (m', n', c_0)^{1/2}c_0^{1/2 + \varepsilon}.
\end{equation*}
For $1 \leq i \leq t$ we write $(m', n', p_i^{s_i}) = p_i^{u_i}$ as well as $m' = m_i' p_i^{u_i}$, $n' = n_i' p_i^{u_i}$. If $s_i < 2r_i + u_i$, we can estimate the sum $S^{\tilde{\gamma}}_i(m'\overline{R_i}, n' \overline{R_i}; c; p_i^{s_i})$ trivially (our bounds are allowed to depend on $q$). Thus, we can assume that $s_i \geq 2r_i + u_i$ for the remainder of this proof. Under this assumption we have $p_i \mid c$, so that $w := 1 + b(\widetilde{\gamma})c$ is a unit modulo $p_i^{s_i}$, and both $a(\widetilde{\gamma})$ and $d(\widetilde{\gamma})$ are units modulo $p_i^{r_i}$. With this notation we can rewrite 
\begin{equation*}
    S^{\tilde{\gamma}}_i(m'\overline{R_i}, n' \overline{R_i}; c; p_i^{s_i}) = \sum_{\substack{a \modulo p_i^{s_i}\\
    a = a(\widetilde{\gamma}) \modulo p_i^{r_i}}} e\left(\frac{m_i' \overline{R_i} a + wn_i' \overline{R_i} \overline{a}}{p_i^{s_i - u_i}}\right) = p_i^{u_i} \sum_{\substack{a \modulo p_i^{s_i - u_i}\\
    a = a(\widetilde{\gamma}) \modulo p_i^{r_i}}} e\left(\frac{m_i' \overline{R_i} a + wn_i' \overline{R_i} \overline{a}}{p_i^{s_i - u_i}}\right).
\end{equation*}
If $v_i = \lceil (s_i - u_i)/2\rceil$ and $a_1 = a(1 + xp^{v_i})$, then its inverse modulo $p_i^{s_i - u_i}$ is given by $\overline{a_1} = \overline{a}(1 - xp^{v_i})$. It follows that 
\begin{equation}
\begin{aligned}
    S^{\tilde{\gamma}}_i(m'\overline{R_i}, n' \overline{R_i}; c; p_i^{s_i}) & = p^{u_i} \sum_{\substack{a \modulo p_i^{v_i}\\
    a = a(\widetilde{\gamma}) \modulo p_i^{r_i}}} e\left(\frac{m_i' \overline{R_i} a + wn_i' \overline{R_i} \overline{a}}{p_i^{s_i - u_i}}\right)\sum_{x \modulo p^{s_i - u_i - v_i}} e\left(\frac{x(m_i'\overline{R_i}a - wn_i'\overline{R_i}\overline{a})}{p_i^{s_i - u_i - v_i}}\right)\\
    & \leq p^{s_i - v_i} \#\{a \modulo p_i^{v_i} : a^2 m_i' = wn_i' \modulo p^{s_i - u_i - v_i}\}\\
    & \ll p_i^{v_i + u_i},
\end{aligned}
\end{equation}
where we have used Hensel's lemma in the last step. Putting together the bounds for $i = 0$ and $1 \leq i \leq t$ completes the proof of the lemma.
\end{proof}

\end{document}
